\section{The \thunderbolt Protocol}\label{sec:construction}

This section presents the \thunderbolt protocol in full. We begin with the design rationale (\autoref{sec:design-rationale}), explaining the role of the Thunderbolt Ledger, the choice of VEDL for cross-domain binding, and the necessity of the VEDL+VESS composition. We then fix notation and describe the complete protocol flow (\autoref{sec:lifecycle}), followed by the two core zero-knowledge proofs that enforce per-step correctness: the Invoice Consistency Proof (\autoref{sec:proof1}) and the Transfer Consistency Proof (\autoref{sec:proof2}). Their joint consequence (\autoref{sec:joint}) establishes the fixed-key invariant. Finally, we show how to distribute the semi-trusted party into a threshold committee via VESS (\autoref{sec:vess-committee}).

\subsection{Design Rationale}\label{sec:design-rationale}

The introduction motivates the key-splitting insight and the two-phase transfer flow. This subsection explains the three design decisions that make the construction concrete: why the protocol needs a persistent ledger with per-channel binding, why VEDL is the right cross-domain proof technique, and why neither VEDL nor VESS alone suffices.

\emph{The Thunderbolt Ledger and context identifiers.} Because neither the sender nor the recipient need be online at the same time, the protocol requires a persistent medium where invoices and credentials can be deposited and later retrieved. The \emph{Thunderbolt Ledger} (\autoref{sec:sm-participants}) serves this role. Each UTXO channel is bound to the ledger via a distinct \emph{context identifier} $\mathsf{ctx\_id}$ that enters every Fiat--Shamir challenge hash (e.g., $c = \Hash(\ldots, \mathsf{ctx\_id}, \ldots)$), cryptographically binding each proof to its specific channel and preventing cross-channel replay.

\emph{VEDL for cross-domain binding.} The core challenge is binding the recipient's refresh value $r_s$ consistently across three distinct algebraic structures: an elliptic-curve point and two Paillier ciphertexts. A generic Sigma protocol could prove each relation independently, but independent proofs would not prevent the prover from using different values in different domains. Verifiable encryption of discrete logarithms~\cite{camenisch2003verifiable} (\autoref{sec:bg-vedl}) solves this with a single shared Sigma response $z$ that must simultaneously satisfy both an EC and a Paillier verification equation, forcing the same integer into both. We instantiate VEDL twice: the Invoice Consistency Proof extends it to a multi-receiver setting (one EC commitment and two Paillier ciphertexts bound by a single $z$); the Transfer Consistency Proof uses a standard single-receiver instance.

\emph{VESS for threshold trust.} VEDL alone concentrates the semi-trusted party role in a single actor. VESS~\cite{shoup2024simple} (\autoref{sec:bg-vess}) removes this single point of failure: $r_s$ is verifiably encrypted as a $(t,n)$-Shamir sharing, so that any $t$ honest members can perform $y_{s+1} = y_s - r_s$ collectively, while no coalition of fewer than $t$ learns anything. This substitution is modular: VESS replaces the single ciphertext $\mathrm{ct}_C$ with a set of verifiably encrypted shares; the Transfer Consistency Proof is unchanged, and $\pi_{\mathrm{inv}}$ narrows to the two-receiver relation $(D_s, \mathrm{ct}_B)$ while a separate VESS consistency proof covers the committee shares.

\emph{Necessity of the composition.} Neither primitive alone suffices. VEDL without VESS forces a single trusted third party. VESS without VEDL cannot bind the shared secret to the on-chain elliptic-curve anchor. Together, VEDL enforces per-step correctness and VESS distributes trust across a threshold group.

\subsection{Protocol Description}\label{sec:lifecycle}

Before formalizing the consistency proofs, we describe the complete protocol flow. A \thunderbolt channel passes through three phases: \textbf{Open}, \textbf{Transfer} (zero or more), and \textbf{Finalize}---illustrated end-to-end in \autoref{fig:overview}. At every step $s$, the protocol maintains a \emph{holder secret} $x_s$, a \emph{committee secret} $y_s$, and a \emph{public anchor} $X_s = x_s\,H$, subject to the fixed-key invariant
\begin{equation}\label{eq:invariant}
P = (x_s + y_s)\,G.
\end{equation}

\paragraph{Notation.} Throughout, we use the notation of \autoref{sec:bg-ec}: $\GG$ is secp256k1 with prime order $q$, standard base point $G$ (used for the lock key $P$), and NUMS generator $H$ (used for all anchors and deltas). Let $(N_B, g_B)$ and $(N_C, g_C)$ be Paillier public keys for Bob and the Committee, respectively. All Paillier arithmetic is modulo $N_B^2$ or $N_C^2$ as appropriate. All elliptic-curve scalar equalities are taken modulo $q$. Let $\kappa$ denote the statistical security parameter (concretely, $\kappa = 128$); we set $\tilde{q} = q^2 \cdot 2^{\kappa}$ as the masking range for Sigma-protocol commitment scalars, which ensures statistical zero-knowledge with distance at most $2^{-\kappa}$.

\paragraph{Running example.} To ground the notation: suppose at step $s{=}2$ Alice holds secret $x_2$ with public anchor $X_2 = x_2\,H$, and the committee holds $y_2$ such that $(x_2 + y_2)\,G = P$. Bob picks a fresh $r_2$, publishes three invoice objects $(D_2 = r_2\,H,\; \mathrm{ct}_B = \Enc_{pk_B}(r_2),\; \mathrm{ct}_C = \Enc_{pk_C}(r_2))$ plus proof $\pi_{\mathrm{inv}}$. Alice folds: $T_3 = \Enc_{pk_B}(x_2) \oplus \mathrm{ct}_B = \Enc_{pk_B}(x_2 + r_2)$, updates $X_3 = X_2 + D_2$, and publishes $\pi_{\mathrm{xfer}}$. Bob decrypts $T_3$ to obtain $x_3 = x_2 + r_2$; the committee computes $y_3 = y_2 - r_2$. The lock key $P$ never changes.

\paragraph{Open (on-chain).}
The initial holder $\mathsf{H}_0$ and the Committee jointly derive $P = (x_0 + y_0)\,G$ and fund a Taproot UTXO~\cite{bip340} locked to $P$; in the threshold setting, at least $t$ committee members must acknowledge the opening. The initial anchor is $X_0 = x_0\,H$ and the initial holder ciphertext is $T_0 = \Enc_{pk_{\mathsf{H}_0}}(x_0)$. This corresponds to step~\ding{172} in \autoref{fig:overview}.

\paragraph{Transfer (off-chain).}

\input{figures/fig-proof-arch}

Each ownership transition is a purely off-chain operation mediated by the Thunderbolt Ledger. Denote the current holder by Alice and the next recipient by Bob. A single transfer from Alice to Bob proceeds in two phases (steps~\ding{173}--\ding{175} in \autoref{fig:overview}):

\begin{enumerate}
  \item \textbf{Invoice phase (Bob).} Bob chooses a fresh refresh value $r_s \in \ZZ_q$ and publishes the following to the Thunderbolt Ledger:
  \begin{itemize}
    \item a public delta point $D_s = r_s\,H$,
    \item a ciphertext to himself, $\mathrm{ct}_B = \Enc_{pk_B}(r_s;\, \rho_B)$,
    \item a ciphertext to the Committee, $\mathrm{ct}_C = \Enc_{pk_C}(r_s;\, \rho_C)$,
    \item a proof $\pi_{\mathrm{inv}}$ (the Invoice Consistency Proof) that all three objects encode the same $r_s$.
  \end{itemize}

  \item \textbf{Transfer phase (Alice).} Alice fetches the invoice from the Thunderbolt Ledger, verifies $\pi_{\mathrm{inv}}$, forms the next holder ciphertext
  $$
  T_{s+1} = \Enc_{pk_B}(x_s;\, \eta_s) \oplus \mathrm{ct}_B,
  $$
  updates the public anchor to $X_{s+1} = X_s + D_s$, and publishes $(T_{s+1}, X_{s+1}, \pi_{\mathrm{xfer}})$ to the Thunderbolt Ledger.
\end{enumerate}

After the transfer, Bob fetches $T_{s+1}$ from the ledger, decrypts it to obtain $x_s + r_s$, and sets $x_{s+1} \equiv x_s + r_s \pmod{q}$. The Committee fetches and verifies both $\pi_{\mathrm{inv}}$ and $\pi_{\mathrm{xfer}}$ against the current anchor $X_s$, then recovers $r_s$ from $\mathrm{ct}_C$ and updates $y_{s+1} = y_s - r_s$. No step requires Alice and Bob to be online simultaneously or touches the chain. The fixed-key invariant~\eqref{eq:invariant} is preserved:
$$
(x_{s+1} + y_{s+1})\,G = (x_s + r_s + y_s - r_s)\,G = (x_s + y_s)\,G = P.
$$

\paragraph{Finalize (on-chain).} To close the channel, the holder first proves knowledge of the current holder secret $x_s$ by presenting a Schnorr proof of knowledge of the discrete log of $X_s$ with respect to $H$. Each committee member verifies this proof against the current anchor before agreeing to participate. The holder and the Committee then jointly sign the closing transaction under $P$ via a standard two-party Schnorr signing protocol~\cite{lindell2017fast} (step~\ding{176} in \autoref{fig:overview}); in the threshold setting, at least $t$ committee members must participate in the signing. The chain sees exactly two transactions regardless of how many transfers occurred.

By \autoref{thm:correct-step} and induction on $s$, the protocol preserves three invariants at every step: $(x_s + y_s)\,G = P$ (fixed key), $X_s = x_s\,H$ (anchor correctness), and $T_s = \Enc_{pk_{\mathsf{H}_s}}(x_s)$ (ciphertext correctness).

\paragraph{Security assumptions.}
Both proofs rely on ECDLP hardness in $\GG$, the DCR assumption~\cite{paillier1999public} for the Paillier moduli, and the random-oracle model~\cite{bellare1993random} for Fiat--Shamir. Under these, both are complete, specially sound, and honest-verifier zero-knowledge (HVZK) as Sigma protocols, and NIZK after Fiat--Shamir; soundness follows from the forking lemma~\cite{pointcheval1996security}.

\subsection{Invoice Consistency Proof}\label{sec:proof1}

The invoice phase requires Bob to publish three objects encoding the same refresh value $r_s$: an elliptic-curve commitment $D_s$, a Paillier ciphertext $\mathrm{ct}_B$ for himself, and a Paillier ciphertext $\mathrm{ct}_C$ for the Committee. The Invoice Consistency Proof $\pi_{\mathrm{inv}}$ is a NIZK that forces all three to encode the same integer, preventing Bob from breaking the fixed-key invariant by committing inconsistent values.

\begin{definition}[Invoice Consistency Relation]\label{def:inv-relation}
The relation $\mathcal{R}_{\mathrm{inv}}$ is defined as
$$
\mathcal{R}_{\mathrm{inv}} = \left\{
  (D_s, \mathrm{ct}_B, \mathrm{ct}_C ;\; r_s, \rho_B, \rho_C) \;:\;
  \begin{array}{l}
  D_s = r_s\,H, \\[2pt]
  \mathrm{ct}_B = g_B^{r_s}\, \rho_B^{N_B} \bmod N_B^2, \\[2pt]
  \mathrm{ct}_C = g_C^{r_s}\, \rho_C^{N_C} \bmod N_C^2
  \end{array}
\right\}.
$$
\end{definition}

\paragraph{Sigma protocol.} The prover (Bob) samples a commitment scalar $\alpha \gets \ZZ_{\tilde{q}}$ and commitment randomnesses $\beta_B \gets \ZZ_{N_B}^*$, $\beta_C \gets \ZZ_{N_C}^*$, and computes:
\begin{align*}
A_D &= \alpha\,H, \\
A_B &= g_B^{\alpha}\, \beta_B^{N_B} \bmod N_B^2, \\
A_C &= g_C^{\alpha}\, \beta_C^{N_C} \bmod N_C^2.
\end{align*}

\paragraph{Fiat--Shamir challenge.} The challenge is computed as
\begin{multline*}
c = \Hash\bigl(\domaintag{tb/invoice},\; \domaintag{ctx\_id},\; \domaintag{inv\_id},\; H,\; pk_B,\; pk_C, \\
D_s,\; \mathrm{ct}_B,\; \mathrm{ct}_C,\; A_D,\; A_B,\; A_C\bigr) \bmod q.
\end{multline*}

\paragraph{Response.} The prover computes a single shared response and two Paillier randomness responses:
\begin{align*}
z &= \alpha + c \cdot r_s, \\
\tau_B &= \beta_B \cdot \rho_B^c \bmod N_B, \\
\tau_C &= \beta_C \cdot \rho_C^c \bmod N_C.
\end{align*}

The proof is $\pi_{\mathrm{inv}} = (A_D, A_B, A_C, z, \tau_B, \tau_C)$.

\paragraph{Verification.} The verifier accepts iff all three of the following hold (same $z$ in all three):
\begin{enumerate}
  \item \textbf{EC check:}
  $
  z\,H \stackrel{?}{=} A_D + c \cdot D_s.
  $
  \item \textbf{Bob--Paillier check:}
  $
  g_B^{z}\, \tau_B^{N_B} \stackrel{?}{=} A_B \cdot \mathrm{ct}_B^{c} \pmod{N_B^2}.
  $
  \item \textbf{Committee--Paillier check:}
  $
  g_C^{z}\, \tau_C^{N_C} \stackrel{?}{=} A_C \cdot \mathrm{ct}_C^{c} \pmod{N_C^2}.
  $
\end{enumerate}

\begin{theorem}[Soundness of the Invoice Consistency Proof]\label{thm:proof1-sound}
The Invoice Consistency Proof is a NIZK for $\mathcal{R}_{\mathrm{inv}}$ in the random-oracle model. In particular, the shared response $z$ in all three verification equations forces a single integer $r_s$ to be the discrete log of $D_s$ and the Paillier plaintext of both $\mathrm{ct}_B$ and $\mathrm{ct}_C$.
\end{theorem}

The complete proof (completeness, special soundness, and honest-verifier zero-knowledge) is given in Appendix~\ref{app:proof-inv}.


\subsection{Transfer Consistency Proof}\label{sec:proof2}

In the transfer phase, Alice folds her current holder secret $x_s$ into the recipient's ciphertext $\mathrm{ct}_B$ to produce the next holder ciphertext $T_{s+1}$. The Transfer Consistency Proof $\pi_{\mathrm{xfer}}$ is a NIZK that forces the folded value to match the public anchor $X_s$, preventing Alice from injecting an arbitrary secret that would leave Bob with an unusable credential. This proof assumes the Invoice Consistency Proof has already been verified, so $D_s$ and $\mathrm{ct}_B$ are known to be consistent with a single $r_s$.

\begin{definition}[Transfer Consistency Relation]\label{def:xfer-relation}
The relation $\mathcal{R}_{\mathrm{xfer}}$ is defined as
$$
\mathcal{R}_{\mathrm{xfer}} = \left\{
  (X_s, \mathrm{ct}_B, T_{s+1} ;\; x_s, \eta_s) \;:\;
  \begin{array}{l}
  X_s = x_s\,H, \\[2pt]
  T_{s+1} \cdot \mathrm{ct}_B^{-1} = g_B^{x_s}\, \eta_s^{N_B} \bmod N_B^2
  \end{array}
\right\}.
$$
\end{definition}

\paragraph{Sigma protocol.} The prover (Alice) samples a commitment scalar $\alpha \gets \ZZ_{\tilde{q}}$ and a Paillier commitment randomness $\beta \gets \ZZ_{N_B}^*$, and computes:
\begin{align*}
A_X &= \alpha\,H, \\
A_T &= g_B^{\alpha}\, \beta^{N_B} \bmod N_B^2.
\end{align*}

\paragraph{Fiat--Shamir challenge.}
\begin{multline*}
c = \Hash\bigl(\domaintag{tb/xfer},\; \domaintag{ctx\_id},\; s{+}1,\; \domaintag{inv\_id},\; H, \\
pk_B,\; X_s,\; \mathrm{ct}_B,\; T_{s+1},\; A_X,\; A_T\bigr) \bmod q.
\end{multline*}

\paragraph{Response.}
\begin{align*}
z &= \alpha + c \cdot x_s, \\
\tau &= \beta \cdot \eta_s^c \bmod N_B.
\end{align*}

The proof is $\pi_{\mathrm{xfer}} = (A_X, A_T, z, \tau)$.

\paragraph{Verification.} The verifier accepts iff both of the following hold (same $z$):
\begin{enumerate}
  \item \textbf{EC check:}
  $
  z\,H \stackrel{?}{=} A_X + c \cdot X_s.
  $
  \item \textbf{Paillier check:}
  $
  g_B^{z}\, \tau^{N_B} \stackrel{?}{=} A_T \cdot (T_{s+1} \ominus \mathrm{ct}_B)^{c} \pmod{N_B^2}.
  $
\end{enumerate}

Additionally, the verifier checks the public anchor transition:
$$
X_{s+1} \stackrel{?}{=} X_s + D_s.
$$
This check is outside the zero-knowledge relation itself but is necessary for protocol correctness.

\begin{theorem}[Soundness of the Transfer Consistency Proof]\label{thm:proof2-sound}
The Transfer Consistency Proof is a NIZK for $\mathcal{R}_{\mathrm{xfer}}$ in the random-oracle model. The shared response $z$ forces a single integer $x_s$ to be both the discrete log of $X_s$ and the Paillier plaintext of $T_{s+1} \ominus \mathrm{ct}_B$.
\end{theorem}

The complete proof follows the same structure as \autoref{thm:proof1-sound} and is given in Appendix~\ref{app:proof-xfer}.


\subsection{Joint Consequence: Correct Transfer Step}\label{sec:joint}

The two proofs above each bind a single witness to one algebraic domain; together, they guarantee that every transfer step preserves the fixed-key invariant. The following theorem makes this precise.

\begin{theorem}[Transfer Correctness]\label{thm:correct-step}
If both $\pi_{\mathrm{inv}}$ and $\pi_{\mathrm{xfer}}$ verify, and the anchor transition $X_{s+1} = X_s + D_s$ holds, then there exist integers $x_s$ and $r_s$ such that
\begin{align}
X_s &= x_s\,H, \label{eq:anchor-open} \\
D_s &= r_s\,H, \label{eq:delta-open} \\
T_{s+1} &= \Enc_{pk_B}(x_s + r_s), \label{eq:next-ct} \\
X_{s+1} &= (x_s + r_s)\,H. \label{eq:next-anchor}
\end{align}
Furthermore, the committee-side update $y_{s+1} = y_s - r_s$ preserves the invariant~\eqref{eq:invariant}:
$$
(x_{s+1} + y_{s+1})\,G = (x_s + y_s)\,G = P.
$$
\end{theorem}

\begin{proof}
From the Invoice Consistency Proof (\autoref{thm:proof1-sound}), extract $r_s$ satisfying $D_s = r_s\,H$ and $\mathrm{ct}_B = \Enc_{pk_B}(r_s)$. From the Transfer Consistency Proof (\autoref{thm:proof2-sound}), extract $x_s$ satisfying $X_s = x_s\,H$ and $T_{s+1} \ominus \mathrm{ct}_B = \Enc_{pk_B}(x_s)$. By Paillier's additive homomorphism, $T_{s+1} = \Enc_{pk_B}(x_s) \oplus \Enc_{pk_B}(r_s) = \Enc_{pk_B}(x_s + r_s)$, giving~\eqref{eq:next-ct}. The anchor transition gives $X_{s+1} = X_s + D_s = x_s\,H + r_s\,H = (x_s + r_s)\,H$, giving~\eqref{eq:next-anchor}. Setting $x_{s+1} \equiv x_s + r_s \pmod{q}$ and $y_{s+1} \equiv y_s - r_s \pmod{q}$, we obtain $(x_{s+1} + y_{s+1})\,G = (x_s + y_s)\,G = P$.
\end{proof}

\subsection{Threshold Committee via Simple VESS}\label{sec:vess-committee}

The construction so far assumes a single third party $\mathsf{C}$. A compromised $\mathsf{C}$ can collude with any holder to recover $x_s + y_s$ and steal the UTXO. We replace $\mathsf{C}$ with a committee of $n$ members via Simple VESS~\cite{shoup2024simple}: at channel open, $y_0$ is distributed as a $(t,n)$-Shamir sharing with Feldman commitments~\cite{feldman1987practical}; at each transfer, Bob replaces $\mathrm{ct}_C$ with a VESS distribution of $r_s$ and attaches a consistency proof. This is a \emph{modular substitution}: the Transfer Consistency Proof remains unchanged, and $\pi_{\mathrm{inv}}$ narrows to the two-receiver relation $(D_s, \mathrm{ct}_B)$ while the committee-side binding is handled by the VESS consistency proof.

\paragraph{Proof bridging.}
In the threshold setting, $\pi_{\mathrm{inv}}$ covers only the two-receiver relation $(D_s, \mathrm{ct}_B)$. The committee-side binding comes from the VESS consistency proof, which guarantees that the encrypted shares encode a unique polynomial whose free-term commitment is $F_0$. The bridge is the public identity $F_0 = D_s$: $\pi_{\mathrm{inv}}$ certifies that $D_s$ commits to $r_s$, and VESS certifies that the shares are consistent evaluations of the same $r_s$.

\paragraph{Committee update and finalization.}
Each member $\mathsf{C}_i$ decrypts its share $r_s^{(i)}$ and updates locally: $y_{s+1}^{(i)} = y_s^{(i)} - r_s^{(i)}$. Since both the committee polynomial and the sharing polynomial for $r_s$ have degree $t{-}1$, their difference is also degree $t{-}1$ with free term $y_{s+1} = y_s - r_s$, preserving the Shamir structure and the fixed-key invariant. At finalization, any $t$ members reconstruct $y_s$ via Lagrange interpolation and jointly sign with the holder via the two-party Schnorr protocol of \autoref{sec:bg-schnorr}.

\paragraph{Security.}
As long as fewer than $t$ members are corrupted, the adversary learns nothing about $y_s$ or $r_s$ (by information-theoretic security of Shamir sharing), and VESS verification prevents a malicious Bob from distributing inconsistent shares. The resulting trust model is the standard $(t,n)$-threshold honesty assumption, identical to FROST~\cite{komlo2021frost} and ROAST~\cite{ruffing2022roast}.

\begin{theorem}[VESS Modularity]\label{thm:vess-modular}
The threshold-committee instantiation preserves every property of the single-third-party construction (fixed-key invariance, anchor correctness, ciphertext binding, and proof soundness) under the $(t,n)$-threshold honesty assumption and the security of Simple VESS~\cite{shoup2024simple}.
\end{theorem}

The complete proof is given in Appendix~\ref{app:proof-vess}.
